Los factores estructurales y contextuales que tienen incidencia en los puntajes de las pruebas Saber 11 han sido estudiados en la literatura colombiana, lo que permite fundamentar qué información se requiere capturar y a partir de está, generar un sistema de información que busque realizar un análisis comparativo y predictivo con robustez.
Uno de los estudios más citados es el que realizó~\cite{chica2010determinantes}, en el año 2010, donde mediante un modelo Logit Ordenado generalizado se encontraron impactos positivos y significativos de variables como nivel de ingreso del hogar, escolaridad de los padres, tipo de jornada escolar y naturaleza del colegio (oficial vs.\ no oficial) sobre los puntajes en matemáticas y lenguaje. El estudio de~\cite{arria2024estudio}, en el año 2024, encontró que los puntajes bajos se asocian sistemáticamente con carencia de recursos educativos, bajo nivel educativo de los padres y condiciones de vida menos favorables, mientras que los puntajes altos se relacionan con mayores ingresos y mejores entornos familiares.
En el Atlántico, la investigación realizada por~\cite{lozanofactores}, en el año 2022, aplicó análisis demográficos, sociales y económicos para identificar variables determinantes del rendimiento en Saber 11, tales como el acceso a internet, características del colegio y variables socioeconómicas del estudiante. Otro estudio regional, realizado por~\cite{miranda2018ingreso}, en el año 2018, encontró que el ingreso del hogar, el estrato socioeconómico y la educación de los padres tienen efectos estadísticamente significativos sobre el puntaje en Saber 11, y que la educación de los padres contribuye de forma similar y consistente al rendimiento académico del hijo.
Adicionalmente, el informe nacional del ICFES del año 2023,~\cite{icfes2024_informeSaber11_2023}, presenta análisis de efectos marginales en modelos estadísticos que muestran cómo factores económicos, familiares y escolares inciden en los resultados del examen Saber 11 a nivel nacional, incluyendo variables como la condición de población víctima del conflicto armado, lo que permite dimensionar desigualdades territoriales y sociales.
Por su parte,~\cite{sanchez2011etnia} en el año 2011, documenta que los estudiantes pertenecientes a etnias cómo indígena, afrocolombiana, raizal y rom, obtienen en promedio puntajes inferiores, y que aunque parte de esa brecha se explica por características observables (familiares, institucionales y sociales), otra parte se debe a factores no observados, lo que implica que el sistema debe contemplar variables latentes o de contexto para capturar esta variación.
Finalmente, el estudio realizado por~\cite{garcia2012factores}, en el año 2012, estudia múltiples grupos de variables independientes ---demográficas, socioeconómicas, familiares, escolares e institucionales--- y encuentra que la zona y condicion de trabajo, pertenecer a una etnia, la pérdida de grados, el tipo de colegio, la jornada escolar y el acceso a internet son determinantes significativos del desempeño destacado frente al no destacado.
Estos antecedentes muestran que los factores que consistentemente aparecen asociados al rendimiento en Saber 11 incluyen condiciones como el nivel socioeconómico, ingresos económicos del hogar, escolaridad de los padres, tipo de colegio (oficial/privado), jornada escolar, acceso a recursos como internet o computadores, etnia, pérdida de grados, así como diferencias regionales y contextuales.
Para el sistema propuesto, esto implica que debe diseñarse para recolectar, limpiar e integrar estas variables, asegurar que los datos sean representativos (territorialidad, diversidad de programas, etc.) e incorporar métodos de modelado que permitan comparar subgrupos, estimar efectos marginales y posiblemente identificar variables no observadas o proxies (sustitutas) que reflejen el contexto institucional. De esta manera, se podrán lograr predicciones útiles y sobre todo, crear y diseñar estrategias institucionales orientadas a fortalecer las competencias académicas de los estudiantes.